\chapter{A Probabilistic Outermost-to-Innermost Transformation}
\label{chap:main-chapter}

This chapter presents the main contribution of the thesis: a transformation
that reduces outermost almost-sure termination of a probabilistic term rewrite
system to innermost almost-sure termination of another probabilistic term
rewrite system.  The construction follows the marker-based architecture of
Thiemann's deterministic transformation \cite{thiemann_2009_outermost}.  Its
administrative rules are deterministic, while the probabilities of the source
rules are retained at the unique point at which a source rewrite step is
simulated.

\Cref{sec:probabilistic-transformation} defines the transformed system and
illustrates a complete marker cycle.  In
\cref{sec:forward-simulation}, every outermost source step is simulated by a
finite sequence of innermost target steps, which yields soundness.  The converse
direction is established in \cref{sec:backward-simulation}: annotated dependency
pairs are used to isolate the only possible source of non-converging target
computations, and a projection then turns such computations into outermost
source computations.  The two directions are combined in
\cref{sec:main-equivalence}.

%-----------------------------------------------------------------------
\section{The Probabilistic Transformation}
\label{sec:probabilistic-transformation}
%-----------------------------------------------------------------------

Let \(\P\) be a finite PTRS over a finite signature \(\Sigma\), and write
\(\operatorname{ar}(f)=n\) whenever \(f\in\Sigma_n\).  The
transformation uses three phases.  First, a \(\treduce\)-marker descends from
the root towards a chosen outermost redex.  Second, the corresponding source
rule is applied and the marker changes into \(\tgoup\).  Third, the
\(\tgoup\)-marker returns to the root, where the next simulation cycle begins.
The intermediate \(\tredex_f\)-term tests whether the current source subterm is
a redex.  If it is, innermost rewriting forces the firing rule to be used before
a descent rule can be applied.  The marker therefore cannot move below a source
redex.

\begin{figure}[htbp]
  \centering
  \begin{tikzpicture}[
    >={Latex[length=2mm]},
    phase/.style={
      draw=RoyalBlue!55,
      rounded corners=2pt,
      fill=RoyalBlue!4,
      text width=3.35cm,
      minimum height=1.35cm,
      align=center,
      font=\small
    },
    arr/.style={->,semithick,draw=black!65},
    note/.style={font=\scriptsize,align=center}
  ]
    \node[phase] (search) at (-4.5,0)
      {\textbf{Search}\\[2pt]
       \(\ttop(\treduce(t))\)\\
       descend to an outermost redex};
    \node[phase] (fire) at (0,0)
      {\textbf{Fire}\\[2pt]
       apply one source rule\\
       with unchanged weights};
    \node[phase] (return) at (4.5,0)
      {\textbf{Return}\\[2pt]
       \(\ttop(\tgoup(s))\)\\
       move the marker back to the root};
    \draw[arr] (search) -- node[above,note] {\(\tredex_f\)} (fire);
    \draw[arr] (fire) -- node[above,note] {\(\tresult\)} (return);
    \draw[arr,bend left=28] (return) to node[below,note]
      {reset for the next source step} (search);
  \end{tikzpicture}
  \caption{The three phases of one simulation cycle.}
  \label{fig:transformation-phases}
\end{figure}

We now define the construction formally.  All newly introduced symbols are
fresh, so none of them occurs in \(\Sigma\).

\begin{Definition}{Probabilistic outermost-to-innermost transformation}
  {probabilistic-transformation}
  For every \(f\in\Sigma_n\), introduce the unary symbol \(\tcheck_f\), the
  \(n\)-ary symbol \(\tredex_f\), and, for every \(1\leq i\leq n\), the
  \(n\)-ary symbol \(\tin_{f,i}\).  In addition, introduce the unary symbols
  \(\ttop\), \(\treduce\), \(\tgoup\), and \(\tresult\).  The extended
  signature is
  \[
    \begin{aligned}
      \Sigma'={}&\Sigma
      \uplus\{\ttop,\treduce,\tgoup,\tresult\}\\
      &\uplus\{\tcheck_f,\tredex_f\mid f\in\Sigma\}
      \uplus\{\tin_{f,i}\mid f\in\Sigma_n,\ 1\leq i\leq n\}.
    \end{aligned}
  \]

  The transformed PTRS \(\P'\) over \(\Sigma'\) consists of the following
  rules.  In the schemes, \(f\in\Sigma_n\), and the variables displayed in an
  administrative rule are pairwise distinct.  For \(n=0\), empty argument
  lists and their parentheses are omitted.
  \begin{alignat}{2}
    \treduce(f(x_1,\ldots,x_n))
      &\to\{1:\tcheck_f(\tredex_f(x_1,\ldots,x_n))\}
      && \tag{T1}\label{rule:search}\displaybreak[0]\\
    \tcheck_f(\tredex_f(x_1,\ldots,x_n))
      &\to
      \{1:\tin_{f,i}(x_1,\ldots,\treduce(x_i),\ldots,x_n)\}
      && \tag{T2}\label{rule:descend}\\
    \tredex_f(\ell_1,\ldots,\ell_n)
      &\to\{p_1:\tresult(r_1),\ldots,p_k:\tresult(r_k)\}
      && \tag{T3}\label{rule:fire}\\
    \tcheck_f(\tresult(x))
      &\to\{1:\tgoup(x)\}
      && \tag{T4}\label{rule:start-return}\\
    \tin_{f,i}(x_1,\ldots,\tgoup(x_i),\ldots,x_n)
      &\to\{1:\tgoup(f(x_1,\ldots,x_n))\}
      && \tag{T5}\label{rule:return}\\
    \ttop(\tgoup(x))
      &\to\{1:\ttop(\treduce(x))\}.
      && \tag{T6}\label{rule:reset}
  \end{alignat}
  Here, \eqref{rule:descend} and \eqref{rule:return} are included for every
  \(1\leq i\leq n\).  For every source rule
  \[
    f(\ell_1,\ldots,\ell_n)
      \to\{p_1:r_1,\ldots,p_k:r_k\}\in\P,
  \]
  rule \eqref{rule:fire} is included with exactly the same probabilities.
  Finally, \eqref{rule:reset} is included once.  We write
  \[
    \operatorname{enc}(t)=\ttop(\treduce(t))
  \]
  for the designated encoding of a source term.
\end{Definition}

All rules except \eqref{rule:fire} have singleton right-hand sides.  Thus,
administrative rewriting does not create probabilistic branching.  If several
descent rules or several matching source rules are available, their selection
remains nondeterministic, exactly as rule and redex selection is
nondeterministic in the source system.

The construction is finite and preserves the variable condition of PTRSs.  It
introduces
\[
  4+2|\Sigma|+\sum_{f\in\Sigma}\operatorname{ar}(f)
\]
fresh symbols and
\[
  1+|\P|+2|\Sigma|
    +2\sum_{f\in\Sigma}\operatorname{ar}(f)
\]
rules.  In particular, its size is linear in the source rules and in the total
arity of the source signature when size is measured by the numbers of symbols
and rules.

\Needspace{11\baselineskip}
The following elementary facts will be used repeatedly.

\begin{Lemma}{Passive source terms and root-redex correspondence}
  {passive-source-terms}
  Let \(t,t_1,\ldots,t_n\in\TSet{\Sigma}{\Var}\) and
  \(f\in\Sigma_n\).
  \begin{enumerate}
    \item The term \(t\) is a normal form with respect to \(\P'\).
    \item The term \(\tredex_f(t_1,\ldots,t_n)\) is reducible by \(\P'\) if
          and only if \(f(t_1,\ldots,t_n)\) is a root redex of \(\P\).
  \end{enumerate}
\end{Lemma}

\begin{proof}
  Every left-hand side of \(\P'\) has a fresh administrative root symbol.
  Hence no subterm of a term over \(\Sigma\) can be a \(\P'\)-redex, which
  proves the first claim.

  The only rules with root \(\tredex_f\) are the instances of
  \eqref{rule:fire}.  Such a rule matches
  \(\tredex_f(t_1,\ldots,t_n)\) exactly when the corresponding source
  left-hand side \(f(\ell_1,\ldots,\ell_n)\) matches
  \(f(t_1,\ldots,t_n)\) at the root.
\end{proof}

\begin{example}[Transformation of the running example]
  \label{ex:transformed-producer}
  Consider \(\P_{\mathsf{ex}}\) from
  \cref{ex:probabilistic-strategy-difference}.  Put
  \[
    u_c=\mathsf{first}(\mathsf{pair}(c,\mathsf{produce}))
    \qquad(c\in\{\mathsf a,\mathsf b\})
  \]
  and \(\mu=\{\tfrac12:u_{\mathsf a},\tfrac12:u_{\mathsf b}\}\).  The
  outermost source step at position \(1\) is simulated as follows; all displayed
  arrows are lifted innermost steps of \(\P_{\mathsf{ex}}'\):
  \[
    \begin{aligned}
      \{1:\operatorname{enc}(\mathsf{first}(\mathsf{produce}))\}
      &\iliftto
      \{1:\ttop(\tcheck_{\mathsf{first}}(
        \tredex_{\mathsf{first}}(\mathsf{produce})))\}\\
      &\iliftto
      \{1:\ttop(\tin_{\mathsf{first},1}(
        \treduce(\mathsf{produce})))\}\\
      &\iliftto
      \{1:\ttop(\tin_{\mathsf{first},1}(
        \tcheck_{\mathsf{produce}}(\tredex_{\mathsf{produce}})))\}\\
      &\iliftto
      \left\{
        \begin{array}{l}
          \tfrac12:\ttop(\tin_{\mathsf{first},1}(
            \tcheck_{\mathsf{produce}}(
              \tresult(\mathsf{pair}(\mathsf a,\mathsf{produce}))))),\\[-1pt]
          \tfrac12:\ttop(\tin_{\mathsf{first},1}(
            \tcheck_{\mathsf{produce}}(
              \tresult(\mathsf{pair}(\mathsf b,\mathsf{produce})))))
        \end{array}
      \right\}\\
      &\iliftto^{,*}
      \{\tfrac12:\ttop(\tgoup(u_{\mathsf a})),
        \tfrac12:\ttop(\tgoup(u_{\mathsf b}))\}\\
      &\iliftto
      \operatorname{enc}_{\#}(\mu).
    \end{aligned}
  \]
  The system subscript \(\P'_{\mathsf{ex}}\) has been omitted from the arrows
  to keep the display readable.  The first
  \(\tredex_{\mathsf{first}}\)-term is not reducible, so
  innermost rewriting permits descent into its argument.  At
  \(\tredex_{\mathsf{produce}}\), the transformed probabilistic source rule
  must fire.  The resulting branches have exactly the probabilities of the
  source step.
\end{example}

%-----------------------------------------------------------------------
\section{Forward Simulation and Soundness}
\label{sec:forward-simulation}
%-----------------------------------------------------------------------

For a position \(\pi\), let \(|\pi|\) denote its length.  Thus
\(|\varepsilon|=0\) and \(|i.\pi|=1+|\pi|\).  Let
\(G:\TSet{\Sigma'}{\Var}\to\TSet{\Sigma'}{\Var}\) be defined by
\(G(t)=\tgoup(t)\).  Its pointwise lifting \(G_{\#}\) to
multi-distributions is the lifting from
\cref{def:finite-multidistributions}.

\Needspace{9\baselineskip}
\begin{Lemma}{Simulation of one outermost step}{one-step-simulation}
  Let \(t\in\TSet{\Sigma}{\Var}\), and suppose that
  \(t\oto_{\P,\pi}\mu\), where
  \(\mu=\{p_1:s_1,\ldots,p_k:s_k\}\).  Then
  \[
    \{1:\treduce(t)\}
      \iliftto_{\P'}^{,3(|\pi|+1)}
      G_{\#}(\mu).
  \]
\end{Lemma}

Thus the target derivation is finite and non-empty, and it preserves every
outcome and its probability.

\begin{proof}
  We use induction on \(|\pi|\).

  First let \(\pi=\varepsilon\).  There are a source rule
  \[
    \ell\to\{p_1:r_1,\ldots,p_k:r_k\}\in\P
  \]
  and a substitution \(\sigma\) such that \(t=\ell\sigma\) and
  \(s_j=r_j\sigma\) for every \(j\).  Write
  \(\ell=f(\ell_1,\ldots,\ell_n)\).  By construction of \(\P'\),
  \[
    \begin{aligned}
      \{1:\treduce(t)\}
      &\iliftto_{\P'}
        \{1:\tcheck_f(\tredex_f(\ell_1\sigma,\ldots,
                                  \ell_n\sigma))\}\\
      &\iliftto_{\P'}
        \{p_1:\tcheck_f(\tresult(r_1\sigma)),\ldots,
          p_k:\tcheck_f(\tresult(r_k\sigma))\}\\
      &\iliftto_{\P'}
        \{p_1:\tgoup(r_1\sigma),\ldots,
          p_k:\tgoup(r_k\sigma)\}
       =G_{\#}(\mu).
    \end{aligned}
  \]
  These steps are innermost by
  \cref{lem:passive-source-terms}.  In particular, the reducible
  \(\tredex_f\)-subterm prevents a descent step from being innermost.

  Now let \(\pi=i.\tau\).  Write
  \(t=f(t_1,\ldots,t_i,\ldots,t_n)\).  The contracted source redex lies in
  \(t_i\) at position \(\tau\), so there is a multi-distribution
  \(\nu=\{p_1:u_1,\ldots,p_k:u_k\}\) with
  \(t_i\oto_{\P,\tau}\nu\), and
  \[
    \mu=\{p_j:f(t_1,\ldots,u_j,\ldots,t_n)\mid 1\leq j\leq k\}.
  \]
  The induction hypothesis gives
  \[
    \{1:\treduce(t_i)\}
      \iliftto_{\P'}^{,3(|\tau|+1)}
      \{p_1:\tgoup(u_1),\ldots,p_k:\tgoup(u_k)\}.
  \]
  Since the step at \(i.\tau\) is outermost, \(t\) is not a root redex of
  \(\P\).  By \cref{lem:passive-source-terms},
  \(\tredex_f(t_1,\ldots,t_n)\) is therefore a \(\P'\)-normal form.  Hence
  the descent step is innermost, and the induction hypothesis can be used in
  the surrounding \(\tin_{f,i}\)-context:
  \[
    \begin{aligned}
      \{1:\treduce(t)\}
      &\iliftto_{\P'}
        \{1:\tcheck_f(\tredex_f(t_1,\ldots,t_n))\}\\
      &\iliftto_{\P'}
        \{1:\tin_{f,i}(t_1,\ldots,\treduce(t_i),\ldots,t_n)\}\\
      &\iliftto_{\P'}^{,3(|\tau|+1)}
        \{p_j:\tin_{f,i}(t_1,\ldots,\tgoup(u_j),\ldots,t_n)
          \mid 1\leq j\leq k\}\\
      &\iliftto_{\P'}
        \{p_j:\tgoup(f(t_1,\ldots,u_j,\ldots,t_n))
          \mid 1\leq j\leq k\}\\
      &=G_{\#}(\mu).
    \end{aligned}
  \]
  The last steps are innermost: the unchanged source arguments and the terms
  \(u_j\) are passive, so every proper subterm of the corresponding
  \(\tin_{f,i}\)-redex is a \(\P'\)-normal form.  The number of steps is
  \(2+3(|\tau|+1)+1=3(|\pi|+1)\), as required.
\end{proof}

Adding the surrounding \(\ttop\)-context and one application of
\eqref{rule:reset} gives the encoded form of the simulation.

\begin{Corollary}{Encoded one-step simulation}{encoded-simulation}
  If \(t\oto_{\P,\pi}\mu\), then
  \[
    \{1:\operatorname{enc}(t)\}
      \iliftto_{\P'}^{,3(|\pi|+1)+1}
      \operatorname{enc}_{\#}(\mu).
  \]
\end{Corollary}

\begin{proof}
  Apply \cref{lem:one-step-simulation} below \(\ttop\), and then apply
  \eqref{rule:reset} to every occurrence in the resulting
  multi-distribution.
\end{proof}

The same marker discipline also rules out an infinite computation within a
single traversal.

\begin{Lemma}{Finiteness of one marker traversal}{finite-marker-traversal}
  Let \(t\in\TSet{\Sigma'}{\Var}\) be a \(\P'\)-normal form.  Starting from
  \(\treduce(t)\), every innermost branch either becomes stuck or reaches a
  term \(\tgoup(s)\) after finitely many steps.  The subtree up to this first
  \(\tgoup\)-frontier is finite.
\end{Lemma}

\begin{proof}
  Before a firing rule is used, each descent moves the unique
  \(\treduce\)-marker into a proper subterm of the finite normal form \(t\).
  Hence only finitely many descents are possible.  A branch then either gets
  stuck or uses exactly one firing rule.  In the latter case, every variable
  value used by the matching substitution is a subterm of \(t\), hence a
  normal form.  Since source symbols are passive, the firing outcomes are
  normal forms as well.  Therefore \eqref{rule:start-return} creates
  \(\tgoup\), and every subsequent return
  step removes one stored \(\tin_{f,i}\)-context.  Thus the first
  \(\tgoup\)-frontier is reached after finitely many further steps.  The
  branching is finite and occurs only at the single firing step, so the whole
  first-traversal subtree is finite.
\end{proof}

Source normal forms also produce no administrative divergence.

\begin{Lemma}{Encoded normal forms terminate administratively}
  {encoded-normal-forms}
  If \(t\in\NF_{\oto_{\P}}\), then every innermost derivation from
  \(\operatorname{enc}(t)\) is finite.  In particular, every maximal such
  derivation ends in a \(\P'\)-normal form.
\end{Lemma}

\begin{proof}
  Since all three source strategies have the same normal forms,
  \(t\) contains no \(\P\)-redex.  Consequently, no transformed firing rule
  \eqref{rule:fire} can be used during the search.  Finiteness follows from
  \cref{lem:finite-marker-traversal}.  At a variable, or
  after checking a constant for which no firing rule applies, the marker is
  stuck.  No \(\tgoup\)-marker is created, so neither a return rule nor the
  reset rule can become applicable.
\end{proof}

We can now lift the local simulation to whole rewrite sequence trees.  The RST
formulation is convenient because different nodes may choose outermost redexes
at different depths.  Each source edge is replaced independently by its finite
simulation fragment.

\begin{Theorem}{Soundness}{soundness}
  If \(\P'\) is innermost AST, then \(\P\) is outermost AST:
  \[
    \AST[\ito_{\P'}]\quad\Longrightarrow\quad\AST[\oto_{\P}].
  \]
\end{Theorem}

\begin{proof}
  We prove the contrapositive.  Suppose that \(\P\) is not outermost AST.
  By \cref{def:almost-sure-termination}, there is a non-converging fully
  evaluated outermost \(\P\)-RST \(\mathfrak T\), so
  \(|\mathfrak T|<1\).

  Replace the term component \(t\) of every node label by
  \(\operatorname{enc}(t)\).  At each internal node, unfold the finite lifted
  derivation from \cref{cor:encoded-simulation} into an innermost RST fragment.
  The administrative steps have
  probability one, and \eqref{rule:fire} has exactly the probabilities of the
  corresponding source rule.  The probability of every node is therefore
  unchanged.  At every source leaf, append any maximal administrative
  derivation; it is finite and ends in a \(\P'\)-normal form by
  \cref{lem:encoded-normal-forms}.

  The resulting tree is a fully evaluated innermost \(\P'\)-RST
  \(\mathfrak T'\) with
  \(|\mathfrak T'|=|\mathfrak T|<1\).  Hence \(\P'\) is not innermost AST.
\end{proof}

%-----------------------------------------------------------------------
\section{Backward Simulation and Completeness}
\label{sec:backward-simulation}
%-----------------------------------------------------------------------

The converse direction has to cover every term over \(\Sigma'\), including
terms with several administrative symbols that are not produced by
\(\operatorname{enc}\).  A direct inspection of encoded computations is
therefore insufficient.  Following the role played by ordinary dependency
pairs in Thiemann's completeness proof, we use the complete annotated
dependency-pair (ADP) framework for innermost AST
\cite{kassing_giesl_2026_adp} to extract a non-converging computation of a
canonical shape.

The cited ADP framework assumes that the signature contains a constant.  If
\(\Sigma'\) has none, we adjoin one fresh inert constructor while applying the
framework and continue to write \(\Sigma'\) for the extended signature.  This
does not add a rule or an ADP, and every counterexample RST over the original
signature remains a counterexample after the extension.  The source projection
below maps this target-only constant to a fresh variable.  Thus this harmless
technical extension does not restrict the main theorem.

\subsection{ADP facts used for the extraction}
\label{sec:adp-extraction-facts}

We briefly record the ADP notation needed below.  This is an application of the
existing framework, not part of the transformation itself.  For every defined
symbol \(h\), let \(h^{\sharp}\) be a fresh annotated copy.  If \(r\) is a
term, \(\sharp_{\mathcal D}(r)\) denotes the term obtained by annotating all
  defined symbols in \(r\), and \(\flat(r)\) erases every annotation.  We use
  \(\flat\) pointwise for multi-distributions, ADPs, and sets of ADPs as well.  The
canonical ADP associated with
\(\ell\to\{p_1:r_1,\ldots,p_k:r_k\}\) is
\[
  \ell\to
  \{p_1:\sharp_{\mathcal D}(r_1),\ldots,
    p_k:\sharp_{\mathcal D}(r_k)\}^{\mathsf{true}}.
\]
The set of canonical ADPs of a PTRS \(\mathcal Q\) is denoted by
\(\DP{\mathcal Q}\).  The flag \(\mathsf{true}\) records that the underlying
rule may be used below an annotation while retaining the annotations above the
contracted position.  The non-probabilistic system
\(\operatorname{np}(\mathcal A)\) contains the deannotated support rules of
the true-flag ADPs in \(\mathcal A\).  In particular,
\[
  \NF_{\fto_{\operatorname{np}(\DP{\P'})}}
    =\NF_{\ito_{\P'}}.
\]

An ADP chain tree is an RST over annotated terms in which distinguished
annotated steps track potentially infinite calls; every infinite branch must
contain infinitely many such steps.  Operationally, a step at an annotation
inserts the annotations displayed on the chosen ADP's right-hand side, whereas
a step at an unannotated position uses its deannotated right-hand side.  A
true-flag step below an annotation may retain the annotations above it.  These
are the only details of the ADP rewrite relation needed below.

We use the following established results from
\cite[Thm.~4, Lem.~2, Thms.~8, 10, and~12]{kassing_giesl_2026_adp}:
\begin{enumerate}
  \item the innermost chain criterion states that \(\mathcal Q\) is innermost
        AST if and only if \(\DP{\mathcal Q}\) is innermost AST;
  \item the Starting Lemma permits a non-converging chain tree to be rooted at
        an instantiated, root-annotated left-hand side.  In the innermost case,
        this root is an argument normal form: all its proper subterms are
        normal forms of \(\operatorname{np}(\mathcal A)\);
  \item the innermost dependency-graph, usable-terms, and probabilistic
        reduction-pair processors are sound and complete.
\end{enumerate}
Here soundness means, in particular, that if an input problem is not innermost
AST, then at least one generated subproblem is not innermost AST.  When a chain
tree is turned back into an ordinary RST, annotation erasure preserves every
node probability; this is the probability-preserving completeness construction
in the proof of the chain criterion.

\subsection{Shape of non-converging target computations}
\label{sec:target-shape}

Write \(\mathcal A_0=\DP{\P'}\).  A root-symbol analysis gives the following
sound over-approximation of the relevant part of its innermost dependency
graph.  The only cyclic strongly connected components are contained in
\[
  C_{\mathsf{top}}=\{\rho_{\mathsf{top}}\}
  \qquad\text{and}\qquad
  C_{\mathsf{down}}
   =\{\rho_{\mathsf{red},f},\rho_{\mathsf{desc},f,i}\mid
       f\in\Sigma,\ 1\leq i\leq\operatorname{ar}(f)\},
\]
where the relevant ADPs are represented schematically by
\begin{align*}
  \rho_{\mathsf{top}}:\quad
  \ttop(\tgoup(x))
    &\to\{1:\ttop^{\sharp}(\treduce^{\sharp}(x))\}^{\mathsf{true}},\\
  \rho_{\mathsf{red},f}:\quad
  \treduce(f(x_1,\ldots,x_n))
    &\to\{1:\tcheck_f^{\sharp}(
       \tredex_f^{\sharp}(x_1,\ldots,x_n))\}^{\mathsf{true}},\\
  \rho_{\mathsf{desc},f,i}:\quad
  \tcheck_f(\tredex_f(x_1,\ldots,x_n))
    &\to\{1:\tin_{f,i}^{\sharp}(x_1,\ldots,
       \treduce^{\sharp}(x_i),\ldots,x_n)\}^{\mathsf{true}}.
\end{align*}
Some annotations displayed here are removed by the usable-terms processor.
They are retained in the schematic forms only to show the canonical ADPs before
processing.

Indeed, \(\rho_{\mathsf{top}}\) may lead back to itself or into the downward
region.  A \(\rho_{\mathsf{red},f}\)-ADP may lead to a descent or to an
acyclic firing ADP, while a \(\rho_{\mathsf{desc},f,i}\)-ADP may lead back to
a reduction call or to an acyclic return ADP.  The firing and return families
do not introduce an annotated call that can lead back to them.  These possible
edges yield the schematic condensation in
\cref{fig:transformed-dependency-graph}.

\begin{figure}[htbp]
  \centering
  \begin{tikzpicture}[
    >={Latex[length=2mm]},
    comp/.style={draw,rounded corners=2pt,fill=Gray!8,
                 minimum width=3.3cm,minimum height=1.25cm,
                 align=center,font=\small},
    arr/.style={->,semithick},
    lab/.style={font=\scriptsize,align=center}
  ]
    \node[comp] (top) at (-3.8,0)
      {\(C_{\mathsf{top}}\)\\\(\rho_{\mathsf{top}}\)};
    \node[comp] (down) at (0.5,0)
      {\(C_{\mathsf{down}}\)\\
       \(\rho_{\mathsf{red},f},\rho_{\mathsf{desc},f,i}\)};
    \node[comp] (acyc) at (4.8,0)
      {acyclic firing and\\return ADPs};
    \draw[arr] (top) edge[loop above] node[lab] {reset cycle} (top);
    \draw[arr] (top) -- (down);
    \draw[arr] (down) edge[loop above] node[lab] {descent cycle} (down);
    \draw[arr] (down) -- (acyc);
  \end{tikzpicture}
  \caption{Schematic SCC over-approximation for the innermost dependency graph
    of \(\mathcal A_0\).  Individual downward ADPs that are not cyclic are
    safely included in \(C_{\mathsf{down}}\).}
  \label{fig:transformed-dependency-graph}
\end{figure}

\Needspace{12\baselineskip}
The downward region cannot itself cause non-convergence.  We include the
orientation argument because this is the point at which the probabilistic ADP
framework is essential.  For every cyclic
\(C\subseteq C_{\mathsf{down}}\), the dependency-graph processor produces
\(C\cup\flat(\mathcal A_0\setminus C)\).  After applying the innermost
usable-terms processor, the only right-hand sides that retain annotations have
the following forms:
\[
  \begin{aligned}
    \treduce(f(x_1,\ldots,x_n))
      &\to\{1:\tcheck_f^{\sharp}(
        \tredex_f(x_1,\ldots,x_n))\}^{\mathsf{true}},\\
    \tcheck_f(\tredex_f(x_1,\ldots,x_n))
      &\to\{1:\tin_{f,i}(x_1,\ldots,
        \treduce^{\sharp}(x_i),\ldots,x_n)\}^{\mathsf{true}}.
  \end{aligned}
\]
Consider the interpretation \(A\) given by the following multilinear
polynomials:
\begin{align*}
  A(f(x_1,\ldots,x_n))
    &=1+x_1+\cdots+x_n &&(f\in\Sigma),\\
  A(\tredex_f(x_1,\ldots,x_n))
    &=x_1+\cdots+x_n,\\
  A(\treduce(x))=A(\treduce^{\sharp}(x))
    &=1+x,\\
  A(\tcheck_f(x))=A(\tcheck_f^{\sharp}(x))
    &=1+x,\\
  A(\tin_{f,i}(x_1,\ldots,x_n))
    &=x_i,\\
  A(\tresult(x))=A(\tgoup(x))=A(\ttop(x))
    &=0.
\end{align*}
All remaining annotated symbols, including
\(\tin_{f,i}^{\sharp}\), \(\tredex_f^{\sharp}\), and
\(\ttop^{\sharp}\), are interpreted by zero; the same is done for the fresh
inert constant if it was added.  These polynomials induce a weakly monotone,
\(\IN\)-collapsible barycentric algebra, as required by the probabilistic
reduction-pair processor
\cite[Sec.~4.4]{kassing_giesl_2026_adp}.  All inequalities below are understood
pointwise for every valuation of the variables.  Let
\(A^{\sharp}_{\mathsf{sum}}(r)\) be the sum of the interpretations of all
annotated subterms of \(r\), as in the probabilistic reduction-pair processor.
For every \(\rho_{\mathsf{red},f}\)-ADP we obtain
\[
  \begin{aligned}
    A^{\sharp}_{\mathsf{sum}}(
      \treduce^{\sharp}(f(x_1,\ldots,x_n)))
      &=2+\sum_{j=1}^{n}x_j\\
      &>1+\sum_{j=1}^{n}x_j\\
      &=A^{\sharp}_{\mathsf{sum}}(
        \tcheck_f^{\sharp}(\tredex_f(x_1,\ldots,x_n))).
  \end{aligned}
\]
For every descent ADP,
\[
  1+\sum_{j=1}^{n}x_j
  \geq 1+x_i
  =A^{\sharp}_{\mathsf{sum}}(
    \tin_{f,i}(x_1,\ldots,
      \treduce^{\sharp}(x_i),\ldots,x_n)).
\]
These inequalities verify weak decrease for the two annotated families; all
remaining ADPs have annotation-free right-hand sides and therefore
right-hand-side \(\sharp\)-sum zero.  The strict inequality for
\(\rho_{\mathsf{red},f}\) is condition~(2a) of the processor.  Its additional
deannotated condition~(2b) also holds, since the left-hand side has value
\(2+\sum_jx_j\) and its deannotated right-hand side has value
\(1+\sum_jx_j\).

It remains to check weak expected decrease of the deannotated support rules,
which is condition~(3).  For \eqref{rule:search} and \eqref{rule:descend}, the
two sides have values \(2+\sum_jx_j\geq1+\sum_jx_j\) and
\(1+\sum_jx_j\geq1+x_i\), respectively.  Rules
\eqref{rule:start-return}, \eqref{rule:return}, and \eqref{rule:reset} give
\(1\geq0\), \(0\geq0\), and \(0\geq0\).  The firing rules are the only
non-singleton rules, and they are weakly decreasing in expectation because
\[
  A(\tredex_f(\ell_1,\ldots,\ell_n))
    \geq 0
    =\sum_{j=1}^{k}p_j A(\tresult(r_j)).
\]
Thus all conditions of the probabilistic reduction-pair processor are
satisfied.  It removes the annotations from the
\(\rho_{\mathsf{red},f}\)-ADPs.  A second dependency-graph step then returns no
cyclic subproblem, because every remaining descent annotation can only lead to
one of these now deannotated calls.  Consequently, every subproblem represented
by \(C_{\mathsf{down}}\) is innermost AST.

\begin{Lemma}{Canonical shape of target non-convergence}{canonical-target-shape}
  If \(\P'\) is not innermost AST, then there exists a non-converging
  innermost \(\P'\)-RST \(\mathfrak T\) whose infinite branches repeatedly
  contain fragments of the form
  \[
    \begin{aligned}
      (q_j:\ttop(\tgoup(t_j)))
      &\xrightarrow[\varepsilon]{\mathfrak T}
       (q_j:\ttop(\treduce(t_j)))\\[-1pt]
      &\xrightarrow[>\varepsilon]{\mathfrak T,*}
       (q_{j+1}:\ttop(\tgoup(t_{j+1}))),
    \end{aligned}
    \tag{\(*\)}\label{eq:canonical-fragment}
  \]
  where \(q_j\) and \(q_{j+1}\) are the cumulative probabilities of the
  corresponding nodes and every \(t_j\) is a \(\P'\)-normal form.  The arrows
  in \eqref{eq:canonical-fragment} describe a path in \(\mathfrak T\); they are
  not lifted multi-distribution arrows.  Below every displayed root step, the subtree
  up to the next \(\ttop(\tgoup(\cdot))\)-frontier or an administrative normal
  form is finite.
\end{Lemma}

\begin{proof}
  By the innermost chain criterion,
  \(\mathcal A_0=\DP{\P'}\) is not innermost AST.  Apply the complete
  innermost dependency-graph processor.  All acyclic subproblems are trivially
  innermost AST, and the preceding reduction-pair argument proves innermost AST
  of every downward subproblem.  Therefore the remaining top subproblem
  \[
    \mathcal B_{\mathsf{top}}
      =C_{\mathsf{top}}
        \cup\flat(\mathcal A_0\setminus C_{\mathsf{top}})
  \]
  is not innermost AST.  Applying the innermost usable-terms processor removes
  the non-usable \(\treduce^{\sharp}\)-annotation from the reset ADP, leaving
  \[
    \ttop(\tgoup(x))
      \to\{1:\ttop^{\sharp}(\treduce(x))\}^{\mathsf{true}}.
  \]
  Denote the resulting problem by
  \(\mathcal B_{\mathsf{top}}^{\mathrm{UT}}\).
  Every other ADP in the processed problem has an annotation-free right-hand
  side.  By soundness of the usable-terms processor, this processed problem is
  still not innermost AST.

  The Starting Lemma now gives a non-converging chain tree rooted at an
  argument-normal instance of a left-hand side with only its root annotated.
  This left-hand side must belong to \(\rho_{\mathsf{top}}\).  Otherwise, its
  first annotated step would remove the only annotation, and no resulting
  infinite path could contain the infinitely many annotated steps required of
  a chain tree.  Thus the root has the form
  \(\ttop^{\sharp}(\tgoup(t_1))\).  Moreover, the only annotation that can
  recur is \(\ttop^{\sharp}\).  Consequently, every infinite path repeatedly
  uses the reset ADP at the root; between two such uses, true-flag steps preserve
  the outer annotation while rewriting strictly below it.

  Erase all annotations.  The dependency-graph and usable-terms processors used
  above only erased annotations and retained every \(\mathsf{true}\)-flag.
  Consequently,
  \[
    \operatorname{np}(\mathcal B_{\mathsf{top}}^{\mathrm{UT}})
      =\operatorname{np}(\mathcal A_0),
  \]
  whose normal forms coincide with the \(\P'\)-normal forms by the chain
  criterion.  Every processed ADP deannotates to the corresponding rule of
  \(\P'\), with the same probabilities.
  \Needspace{3\baselineskip}
  Thus the probability-preserving
  erasure construction from the chain criterion preserves innermostness and
  yields a non-converging innermost \(\P'\)-RST.  Its recurring distinguished
  step is precisely the deterministic root step
  \[
    \ttop(\tgoup(t_j))
      \ito_{\P',\varepsilon}
      \{1:\ttop(\treduce(t_j))\}.
  \]
  Between two such steps, rewriting occurs strictly below the persistent
  \(\ttop\)-symbol, which gives \eqref{eq:canonical-fragment}.

  Finally, the displayed root step is innermost.  Hence its proper subterm
  \(\tgoup(t_j)\), and therefore \(t_j\), is a \(\P'\)-normal form.  Applying
  \cref{lem:finite-marker-traversal} below \(\ttop\) proves finiteness of every
  first-return subtree claimed in the statement.
\end{proof}

\subsection{Projection to source terms}
\label{sec:source-projection}

The terms \(t_j\) in \cref{lem:canonical-target-shape} may contain arbitrary
symbols from \(\Sigma'\).  We therefore project them to source terms.  We
harmlessly enlarge the countably infinite variable supply by a disjoint
countable set of fresh variables; variable renaming identifies the enlarged
supply with the one fixed in \cref{def:terms}.

\begin{Definition}{Source projection}{source-projection}
  Let
  \[
    \mathcal U
      =\{u\in\TSet{\Sigma'}{\Var}\setminus\Var
          \mid\rootterm(u)\in\Sigma'\setminus\Sigma\}.
  \]
  Choose a set \(Z=\{z_u\mid u\in\mathcal U\}\), disjoint from \(\Var\),
  whose elements are pairwise distinct, and put
  \(\widehat{\Var}=\Var\uplus Z\).  Define
  \[
    \operatorname{src}:\TSet{\Sigma'}{\Var}
      \longrightarrow \TSet{\Sigma}{\widehat{\Var}}
  \]
  \Needspace{12\baselineskip}
  by
  \[
    \begin{aligned}
      \operatorname{src}(x)&=x &&(x\in\Var),\\
      \operatorname{src}(f(t_1,\ldots,t_n))
        &=f(\operatorname{src}(t_1),\ldots,
             \operatorname{src}(t_n)) &&(f\in\Sigma),\\
      \operatorname{src}(u)&=z_u
        &&(\rootterm(u)\in\Sigma'\setminus\Sigma).
    \end{aligned}
  \]
  For a target substitution
  \(\delta:\Var\to\TSet{\Sigma'}{\Var}\), define the source substitution
  \(\operatorname{src}(\delta)\) on \(\widehat{\Var}\) by
  \[
    x\operatorname{src}(\delta)=\operatorname{src}(x\delta)
      \quad(x\in\Var),
    \qquad
    z\operatorname{src}(\delta)=z
      \quad(z\in Z).
  \]
  The term map is lifted pointwise to multi-distributions as
  \(\operatorname{src}_{\#}\).
\end{Definition}

Fresh variables deliberately hide the internal structure of administrative
subterms.  Choosing a different variable for every such term makes
\(\operatorname{src}\) injective, which is needed to reflect source matching.

\begin{Lemma}{Projection and matching}{projection-matching}
  Let \(\ell\in\TSet{\Sigma}{\Var}\) and
  \(t\in\TSet{\Sigma'}{\Var}\).
  \begin{enumerate}
    \item For every target substitution \(\delta\),
          \(\operatorname{src}(\ell\delta)
            =\ell\operatorname{src}(\delta)\).
    \item If \(\sigma\) is a substitution over \(\widehat{\Var}\) and
          \(\operatorname{src}(t)=\ell\sigma\), then there is a target
          substitution \(\delta\) such that \(t=\ell\delta\).
  \end{enumerate}
\end{Lemma}

\begin{proof}
  The first claim follows by structural induction on \(\ell\).

  For the second claim, observe first that \(\operatorname{src}\) is
  injective: source-rooted terms retain their root and are handled
  homomorphically, whereas every administrative-rooted term is mapped to its
  own fresh variable.  Every value \(x\sigma\) occurring in the equality
  \(\operatorname{src}(t)=\ell\sigma\) is therefore the projection of a
  unique corresponding subterm of \(t\).  Assign that subterm to
  \(x\delta\).  Repeated occurrences of \(x\) give the same subterm by
  injectivity.  A structural induction on \(\ell\) now yields
  \(t=\ell\delta\).
\end{proof}

We can now reverse a complete marker traversal.  The statement is
multi-distribution-valued: this is where the probabilities of the transformed
firing rule are transferred back to the source system.

\begin{Lemma}{Extraction of one outermost step}{one-step-extraction}
  Let \(t\in\TSet{\Sigma'}{\Var}\) be a \(\P'\)-normal form.  If
  \[
    \{1:\treduce(t)\}
      \iliftto_{\P'}^{+}
      G_{\#}(\mu),
  \]
  then
  \[
    \operatorname{src}(t)
      \oto_{\P}
      \operatorname{src}_{\#}(\mu),
  \]
  and every term in \(\Supp(\mu)\) is a \(\P'\)-normal form.
\end{Lemma}

\begin{proof}
  Normal-form retention may add self-loops after the traversal has finished; we
  discard such a suffix and consider the first multi-distribution whose terms all
  have root \(\tgoup\).
  We use induction on the number of applications of the descent rule
  \eqref{rule:descend}.  Since the derivation reaches a term with root
  \(\tgoup\), its first step must use \eqref{rule:search}.  Thus
  \(t=f(t_1,\ldots,t_n)\) for some \(f\in\Sigma_n\), and all \(t_i\) are
  \(\P'\)-normal forms.

  First suppose that the transformed firing rule is used next.  Then there are
  a source rule
  \[
    f(\ell_1,\ldots,\ell_n)
      \to\{p_1:r_1,\ldots,p_k:r_k\}\in\P
  \]
  and a substitution \(\delta\) with \(t_i=\ell_i\delta\).  The only possible
  continuation to a \(\tgoup\)-rooted multi-distribution uses
  \eqref{rule:start-return}, so
  \(\mu=\{p_1:r_1\delta,\ldots,p_k:r_k\delta\}\).  By
  \cref{lem:projection-matching},
  \[
    \operatorname{src}(t)
      =f(\ell_1,\ldots,\ell_n)\operatorname{src}(\delta)
      \oto_{\P,\varepsilon}
      \{p_j:r_j\operatorname{src}(\delta)\mid1\leq j\leq k\}
      =\operatorname{src}_{\#}(\mu).
  \]
  Each \(r_j\delta\) is a \(\P'\)-normal form: its source symbols are passive,
  and every substitution value that occurs in \(r_j\) also occurs in the
  left-hand side and is therefore a subterm of the normal form \(t\).

  Otherwise, the derivation uses \eqref{rule:descend} for some argument
  \(i\).  It first produces
  \[
    \tin_{f,i}(t_1,\ldots,\treduce(t_i),\ldots,t_n).
  \]
  Before the outer \(\tin_{f,i}\)-rule can return, the inner marker must
  complete a shorter traversal
  \[
    \{1:\treduce(t_i)\}
      \iliftto_{\P'}^{+}G_{\#}(\nu).
  \]
  The induction hypothesis yields
  \(\operatorname{src}(t_i)\oto_{\P}
    \operatorname{src}_{\#}(\nu)\), and all outcomes of \(\nu\) are normal
  forms of \(\P'\).  The final return rule puts these outcomes back into the
  \(i\)-th argument of \(f\).  More explicitly, if
  \(\nu=\{q_h:u_h\mid1\leq h\leq m\}\), then
  \[
    \mu=\{q_h:f(t_1,\ldots,u_h,\ldots,t_n)
             \mid1\leq h\leq m\},
  \]
  and pointwise projection gives the corresponding contextual lifting of the
  source step.

  It remains to show that the resulting contextual source step is outermost.
  If \(\operatorname{src}(t)\) were a root redex of \(\P\), the matching
  reflection in \cref{lem:projection-matching} would imply that
  \(t\) matches the corresponding source left-hand side.  Then
  \(\tredex_f(t_1,\ldots,t_n)\) would be reducible below \(\tcheck_f\), so the
  descent rule could not have been applied innermost.  Thus the root is not a
  source redex.  The induction hypothesis already excludes every redex at an
  ancestor position inside \(t_i\), so no source redex is above the contextual
  step.  It is therefore outermost.  Its outcomes are normal forms of \(\P'\), since the
  unchanged arguments and all outcomes of \(\nu\) are normal and \(f\) is
  passive in \(\P'\).
\end{proof}

\begin{Theorem}{Completeness}{completeness}
  If \(\P\) is outermost AST, then \(\P'\) is innermost AST:
  \[
    \AST[\oto_{\P}]\quad\Longrightarrow\quad\AST[\ito_{\P'}].
  \]
\end{Theorem}

\begin{proof}
  Again we prove the contrapositive.  Suppose that \(\P'\) is not innermost
  AST.  By \cref{lem:canonical-target-shape}, there is a non-converging
  innermost \(\P'\)-RST whose infinite branches repeatedly have the form
  \eqref{eq:canonical-fragment}, with every \(t_j\) in normal form.

  First normalize every finite fragment that ends partway through a marker
  traversal: extend it only until it reaches the next
  \(\ttop(\tgoup(\cdot))\)-frontier or an administrative normal form.  This is
  always a finite expansion by \cref{lem:finite-marker-traversal}.  Replacing a
  frontier leaf by such a finite subtree preserves its total weight, so this
  normalization does not change the convergence probability.

  \Needspace{14\baselineskip}
  Now consider a macro-node \((q:\ttop(\tgoup(t)))\).  Its reset child is
  \((q:\ttop(\treduce(t)))\).  If the subsequent administrative search gets
  stuck, leave the corresponding source macro-node unexpanded.  Otherwise,
  collect all nodes at the next \(\ttop(\tgoup(\cdot))\)-frontier.  Dividing
  their weights by \(q\) gives a multi-distribution
  \(G_{\#}(\mu)\), and the finite first-return subtree represents
  \[
    \{1:\treduce(t)\}\iliftto_{\P'}^{+}G_{\#}(\mu).
  \]
  This works at the multi-distribution level, rather than for just one path: before
  the firing rule there is no probabilistic branching, and after it all
  branches perform only the deterministic return rules.

  Contract this whole first-return subtree to one edge, remove the outer
  administrative context, and label the macro-node by
  \(\operatorname{src}(t)\).  At every macro-node with children,
  \cref{lem:one-step-extraction} proves that the new edge is one outermost
  \(\P\)-step and that the relative child weights are exactly its source
  probabilities.  Hence the contracted object is an outermost \(\P\)-RST.

  Contraction changes neither the probability of a macro-node nor the weight of
  any leaf: administrative edges have weight one, and every firing frontier is
  grouped with all of its probabilistic children.  Thus the source RST has the
  same convergence probability, which is smaller than one.  Its variables lie
  in the countably infinite set \(\widehat{\Var}\); applying a bijective variable
  renaming identifies this with the supply \(\Var\) fixed in the preliminaries
  and preserves all rewrite steps.  Therefore \(\P\) is not outermost AST.
\end{proof}

%-----------------------------------------------------------------------
\section{Main Equivalence}
\label{sec:main-equivalence}
%-----------------------------------------------------------------------

Combining \cref{thm:soundness,thm:completeness} yields the desired result.

\begin{Theorem}{Soundness and completeness of the transformation}{main}
  Let \(\P\) be a finite PTRS over a finite signature \(\Sigma\), and let
  \(\P'\) be the PTRS from
  \cref{def:probabilistic-transformation}.  Then
  \[
    \AST[\oto_{\P}]
      \quad\Longleftrightarrow\quad
    \AST[\ito_{\P'}].
  \]
  The right-hand side quantifies over all terms over the extended signature
  \(\Sigma'\), not only terms of the form \(\operatorname{enc}(t)\).
\end{Theorem}

\begin{proof}
  The implication from right to left is
  \cref{thm:soundness}; the implication from left to right is
  \cref{thm:completeness}.
\end{proof}

The theorem concerns AST.  The simulation preserves the probability of every
source outcome, but it inserts a number of administrative steps that depends on
the depth of the selected redex.  Therefore the result does not by itself imply
an analogous equivalence for PAST or SAST, both of which additionally constrain
expected derivation lengths.
